\section{Manifesto Results: Main Application}\label{sec:min-manifesto-results}

The empirical question is whether the tree preserves a real social-science
measurement target. We reuse the Manifesto/Benoit benchmark: $235$ party
platforms with expert-survey means on six seven-point policy dimensions. Each
pipeline chunks a manifesto into contiguous spans, summarizes each span under
a dimension-specific preservation rubric, recursively merges adjacent
summaries, scores the root, and correlates the root score against the expert
mean. The evaluation is root-observed: every sampled platform already has a
full-document target.

\begin{table}[t]
  \centering
  \resizebox{\linewidth}{!}{% Generated table; do not edit by hand.
% Metric: Pearson r (higher better)
\begin{tabular}{lrrrrrrrr}
\toprule
Method & Econ. & Social & Immig. & EU & Env. & Decent. & Macro & \#/6 \\
\midrule
\multicolumn{9}{l}{\textit{Benoit et al. (2025) reference}} \\
Proprietary ensemble, 18 scores (Fig 1) & $+0.870$ & $+0.920$ & $+0.890$ & $+0.910$ & $+0.820$ & $+0.490$ & $+0.817$ & 6/6 \\
Split-expert reference (Table 3) & $+0.880$ & $+0.910$ & $+0.880$ & $+0.950$ & $+0.840$ & $+0.780$ & $+0.873$ & 6/6 \\
LLaMA-3.3-70B (Table 6) & $+0.840$ & $+0.870$ & $+0.860$ & $+0.860$ & $+0.680$ & $+0.400$ & $+0.752$ & 6/6 \\
DeepSeek-V3 (Table 6) & $+0.840$ & $+0.870$ & $+0.890$ & $+0.860$ & $+0.790$ & $+0.450$ & $+0.783$ & 6/6 \\
Gemma-3-27B-IT (Table 6) & $+0.860$ & $+0.860$ & $+0.890$ & $+0.840$ & $+0.860$ & $+0.450$ & $+0.793$ & 6/6 \\
\midrule
\multicolumn{9}{l}{\textit{Ours: per-dim pipeline $\times$ leaf size (Gemma-4-31B-NVFP4, 1 summarizer + 1 scorer per dim)}} \\
leaf = 64 K chars ($\approx$16K tokens) & $+0.916$ & $+0.817$ & $+0.888$ & $+0.924$ & $+0.854$ & $+0.489$ & $+0.815$ & 6/6 \\
leaf = 32 K chars ($\approx$8K tokens) & $+0.929$ & $+0.857$ & $+0.886$ & $+0.916$ & $+0.877$ & $+0.480$ & $+0.824$ & 6/6 \\
leaf = 24 K chars ($\approx$6K tokens) --- full test n$\approx$215 & $+0.892$ & $+0.840$ & $+0.867$ & $+0.896$ & $+0.814$ & $+0.464$ & $+0.796$ & 6/6 \\
leaf = 16 K chars ($\approx$4K tokens) & $+0.925$ & $+0.847$ & $+0.898$ & $+0.909$ & $+0.831$ & $+0.537$ & $+0.824$ & 6/6 \\
leaf =  8 K chars ($\approx$2K tokens) & $+0.939$ & $+0.869$ & $+0.884$ & $+0.901$ & $+0.839$ & $+0.543$ & $+0.829$ & 6/6 \\
\midrule
\multicolumn{9}{l}{\textit{Ours: combined pipeline $\times$ leaf size (one shared summarizer with union rubric, 6 scores)}} \\
leaf = 64 K chars & $+0.910$ & $+0.867$ & $+0.897$ & $+0.910$ & $+0.795$ & $+0.458$ & $+0.806$ & 6/6 \\
leaf = 32 K chars & $+0.917$ & $+0.849$ & $+0.908$ & $+0.889$ & $+0.749$ & $+0.438$ & $+0.792$ & 6/6 \\
leaf = 24 K chars (full test n=229) & $+0.868$ & $+0.850$ & $+0.880$ & $+0.902$ & $+0.809$ & $+0.396$ & $+0.784$ & 6/6 \\
leaf = 16 K chars & $+0.919$ & $+0.851$ & $+0.916$ & $+0.889$ & $+0.808$ & $+0.441$ & $+0.804$ & 6/6 \\
leaf =  8 K chars & $+0.927$ & $+0.874$ & $+0.911$ & $+0.917$ & $+0.801$ & $+0.413$ & $+0.807$ & 6/6 \\
\midrule
\multicolumn{9}{l}{\textit{Ours: tiny-leaf extensions of the chunk sweep (Gemma-4)}} \\
tree, leaf = 4 K chars & $+0.932$ & $+0.886$ & $+0.885$ & $+0.931$ & $+0.838$ & $+0.580$ & $+0.842$ & 6/6 \\
tree, leaf = 2 K chars & $+0.924$ & $+0.867$ & $+0.887$ & $+0.906$ & $+0.793$ & $+0.584$ & $+0.827$ & 6/6 \\
tree, leaf = 1 K chars (stress test, depth $\geq$6) & $+0.933$ & $+0.859$ & $+0.914$ & $+0.925$ & $+0.838$ & $+0.493$ & $+0.827$ & 6/6 \\
\midrule
\multicolumn{9}{l}{\textit{Ours: concat-no-merge (chunks summarized independently, joined, scored --- tests whether the merge step carries signal)}} \\
concat, leaf = 32 K chars & $+0.929$ & $+0.856$ & $+0.897$ & $+0.920$ & $+0.848$ & $+0.571$ & $+0.837$ & 6/6 \\
concat, leaf = 16 K chars (default) & $+0.931$ & $+0.823$ & $+0.909$ & $+0.912$ & $+0.866$ & $+0.523$ & $+0.827$ & 6/6 \\
concat, leaf =  8 K chars & $+0.932$ & $+0.849$ & $+0.882$ & $+0.915$ & $+0.850$ & $+0.594$ & $+0.837$ & 6/6 \\
\midrule
\multicolumn{9}{l}{\textit{Ours: flat baseline (no chunk, no summary; truncate text and score) --- tests whether tree is needed at all}} \\
flat, truncation = 48 K chars & $+0.936$ & $+0.821$ & $+0.926$ & $+0.889$ & $+0.806$ & $+0.587$ & $+0.827$ & 6/6 \\
flat, truncation = 24 K chars (default) & $+0.929$ & $+0.847$ & $+0.889$ & $+0.882$ & $+0.829$ & $+0.542$ & $+0.820$ & 6/6 \\
flat, truncation = 12 K chars & $+0.929$ & $+0.844$ & $+0.938$ & $+0.880$ & $+0.779$ & $+0.564$ & $+0.822$ & 6/6 \\
flat, truncation =  6 K chars & $+0.926$ & $+0.811$ & $+0.967$ & $+0.881$ & $+0.886$ & $+0.243$ & $+0.786$ & 6/6 \\
\midrule
\multicolumn{9}{l}{\textit{Ours: scorer ablations (Gemma-4, Benoit's GPT-4o summaries held fixed)}} \\
1. Zero-shot scorer & $+0.832$ & $+0.886$ & $+0.858$ & $+0.905$ & $+0.668$ & $+0.311$ & $+0.743$ & 6/6 \\
2. Few-shot optimized scorer & $+0.822$ & $+0.892$ & $+0.853$ & $+0.912$ & $+0.627$ & $+0.297$ & $+0.734$ & 6/6 \\
3. Joint scorer baseline (shared across 6 dims) & $+0.837$ & $+0.881$ & $+0.852$ & $+0.904$ & $+0.658$ & $+0.314$ & $+0.741$ & 6/6 \\
4. Joint scorer, few-shot optimized & $+0.817$ & $+0.891$ & $+0.848$ & $+0.908$ & $+0.700$ & $+0.313$ & $+0.746$ & 6/6 \\
5. Joint scorer, reflective prompt-optimized & $+0.832$ & $+0.882$ & $+0.848$ & $+0.879$ & $+0.688$ & $+0.344$ & $+0.746$ & 6/6 \\
\midrule
\multicolumn{9}{l}{\textit{Ours: exact-model replication (Benoit's Gemma-3-27B-IT BF16)}} \\
Gemma-3-27B scorer-only, OUR scoring rubric & $+0.826$ & $+0.864$ & $+0.846$ & $+0.902$ & $+0.573$ & $+0.329$ & $+0.723$ & 6/6 \\
Gemma-3-27B scorer-only, exact Benoit rubric & $+0.814$ & $+0.895$ & $+0.852$ & $+0.852$ & $+0.616$ & $+0.359$ & $+0.731$ & 6/6 \\
Gemma-3-27B, raw-prompt Benoit format (bare-integer response) & $+0.782$ & $+0.901$ & $+0.854$ & $+0.874$ & $+0.590$ & $+0.368$ & $+0.728$ & 6/6 \\
\bottomrule
\end{tabular}
}
  \caption{\textbf{Pearson $r$ against Benoit expert-survey means.}
  The first block reproduces published Benoit reference rows; the remaining
  blocks report C-TreePO and controls using Gemma-4-31B-NVFP4. The main
  headline is the per-dimension tree at $8\mathrm{K}$-character leaves:
  macro $r=0.829$, above the proprietary ensemble's $0.817$ and the matched
  Gemma-3 open-weight replication's $0.793$.}
  \label{tab:min-benoit-headline}
\end{table}

The six-dimension headline is that one open-weight tree pipeline matches the
published frontier-ensemble benchmark. At $8\mathrm{K}$-character leaves, the
per-dimension C-Tree row reaches macro $0.829$ against expert means, compared
with $0.817$ for Benoit's 18-score proprietary ensemble and $0.793$ for their
Gemma-3-27B open-weight replication. The economic dimension reaches $0.939$ in
the same row, above the published economic split-expert agreement reference
of $0.880$.

The stronger methodological point is granularity. Across the character-leaf
sweep from $4\mathrm{K}$ to $64\mathrm{K}$, the per-dimension tree's macro
correlation varies by only $0.027$. Economic policy remains high across the
same range: $0.916$ at $64\mathrm{K}$, $0.939$ at $8\mathrm{K}$, and
$0.932$ at $4\mathrm{K}$. The pattern is root-observed evidence that the
node/root supervision tradeoff is plausible and worth auditing. The
document-level target remains visible as the unit of work shrinks from the
whole manifesto toward leaves small enough to sample and audit. A completed
node-level certificate would add randomized local-law labels over those
partial-document units.

The economic prompt ladder gives the same pattern on a token leaf-size axis.
The label $f^a g^b$ means that the scorer prompt has gone through $a$
scorer-side optimization rounds and the summarizer/merge prompt through $b$
summarizer-side rounds. In the fully populated Benoit-init lane, the best
external Pearson values are $0.885$ at $1024$ tokens, $0.886$ at $2048$,
$0.886$ at $4096$, and $0.879$ at $8192$. The $0.007$-wide band straddles
the economic split-expert reference across an $8\times$ leaf-size range.
Below roughly $1024$ tokens, performance degrades: $0.861$ at $512$ and
$0.830$ at $256$. The drop is the expected failure mode when leaves become
too small to carry the evidence a later merge needs.

\begin{figure}[t]
  \centering
  \includegraphics[width=0.72\linewidth]{manifesto_fg_ladder_headline.pdf}
  \caption{\textbf{Economic prompt-ladder leaf-size robustness.}
  Per-leaf best external Pearson against Benoit expert means for the
  Benoit-init and raw-init lanes. The dash-dot reference is the economic
  split-expert agreement value $r=0.880$. The raw-init lane includes a
  populated pilot cell at leaf $256$, stage $f^1g^1$, with external
  $r=0.909$ and internal-external gap $0.056$; the six-dimension table and
  broader plateau carry the empirical claim. Full grids appear in
  Appendix~\ref{app:min-manifesto}.}
  \label{fig:v4-manifesto-fg-headline}
\end{figure}

The raw-init lane is useful as a pilot. Its leaf-256, $f^1g^1$ cell reaches
$r=0.909$ with internal-external gap $0.056$, the best populated cell in the
current grid. It shows that the method can find a high-agreement cell without
inheriting Benoit's summaries. The more stable empirical claim comes from the
six-dimension table and the broader leaf-size plateau.

The current Manifesto result sits on the root-observed tier. The Manifesto
Project's expert quasi-sentence codings cover $2{,}157$ platforms and about
$2.27$M coded spans, so they provide a path to the next supervision design:
root scores continue to calibrate the full-document target, while span-level
codings can become leaf and merge oracles by aggregating coded quasi-sentences
inside each sampled span. The quasi-sentence route is the planned local-law
supervision and audit layer. Table~\ref{tab:min-benoit-headline} reports the
root-observed correlations.
